\documentclass{article}

\usepackage{fancyhdr}
\usepackage{extramarks}
\usepackage{amsmath}
\usepackage{amsthm}
\usepackage{amsfonts}
\usepackage{tikz}
\usepackage[plain]{algorithm}
\usepackage{algpseudocode}
\usepackage{listings}
\usepackage{xcolor}
\usepackage{float}

\usetikzlibrary{automata,positioning}

\newcommand{\overbar}[1]{\mkern 1.5mu\overline{\mkern-1.5mu#1\mkern-1.5mu}\mkern 1.5mu}

%
% Basic Document Settings
%

\topmargin=-0.45in
\evensidemargin=0in
\oddsidemargin=0in
\textwidth=6.5in
\textheight=9.0in
\headsep=0.25in

\linespread{1.1}

\pagestyle{fancy}
\lhead{Yousef Alaa Awad}
\chead{\hmwkClass\: \hmwkTitle}
\rhead{\firstxmark}
\lfoot{\lastxmark}
\cfoot{\thepage}

\renewcommand\headrulewidth{0.4pt}
\renewcommand\footrulewidth{0.4pt}

\setlength\parindent{0pt}

%
% Create Problem Sections
%

\setcounter{secnumdepth}{0}

\newcommand{\hmwkTitle}{Assignment\ \#4}
\newcommand{\hmwkClass}{CDA 5106}

%
% Title Page
%

\title{
    \vspace{2in}
    \textmd{\textbf{\hmwkClass:\ \hmwkTitle}}\\
    \normalsize\vspace{0.1in}
}

\author{Yousef Alaa Awad}

\begin{document}

\maketitle
\pagebreak

\section{1. Cache Placement and Replacement}
\textbf{Problem Statement:} We are given a processor with a two-level cache hierarchy. The L1 cache is direct-mapped with two lines, where blocks A, B, C, D map to Line 0, and U, V, W, X map to Line 1. The L2 cache is fully associative with four lines. Both utilize the LRU replacement policy. The L2 cache is non-inclusive of the L1 cache.

\subsection{Access Outcomes}
\begin{table}[H]
	\centering
	\renewcommand{\arraystretch}{1.3}
	\begin{tabular}{|c|c|c|}
		\hline
		\textbf{Access} & \textbf{Outcome for L1} & \textbf{Outcome for L2} \\
		\hline
		A               & Miss                    & Miss                    \\
		\hline
		V               & Miss                    & Miss                    \\
		\hline
		U               & MR-V                    & Miss                    \\
		\hline
		B               & MR-A                    & Miss                    \\
		\hline
		A               & MR-B                    & Hit                     \\
		\hline
		D               & MR-A                    & MR-V                    \\
		\hline
		C               & MR-D                    & MR-U                    \\
		\hline
		A               & MR-C                    & Hit                     \\
		\hline
		B               & MR-A                    & Hit                     \\
		\hline
	\end{tabular}
\end{table}

\textit{Note on Access C:} The L2 cache becomes full with blocks A, V, U, B. The LRU accesses determine eviction. When D is accessed, V is evicted from L2. When C is accessed, U is the LRU block in L2 and is evicted. Because the cache is non-inclusive, the eviction of U from L2 does \textbf{not} trigger a back-invalidation of U in the L1 cache. U remains valid in L1 Line 1.

\subsection{Final Cache Contents}
\begin{table}[H]
	\centering
	\renewcommand{\arraystretch}{1.3}
	\begin{tabular}{|c|c|c|}
		\hline
		\textbf{L1 Cache} & \textbf{Which Block?} & \textbf{What State?} \\
		\hline
		Line 0            & B                     & Valid                \\
		\hline
		Line 1            & U                     & Valid                \\
		\hline
		\textbf{L2 Cache} & \textbf{Which Block?} & \textbf{What State?} \\
		\hline
		Line 0            & A                     & Valid                \\
		\hline
		Line 1            & B                     & Valid                \\
		\hline
		Line 2            & C                     & Valid                \\
		\hline
		Line 3            & D                     & Valid                \\
		\hline
	\end{tabular}
\end{table}

\pagebreak

\section{2. Cache Performance}

\subsection{(a) Average Memory Access Time (AMAT)}
Using the standard baseline parameters: $t_{L1} = 2$, $MR_{L1} = 20\%$, $t_{L2} = 15$, $MR_{L2} = 8\%$, and $t_{mem} = 90$.
\begin{align*}
	AMAT & = t_{L1} + MR_{L1} \times (t_{L2} + MR_{L2} \times t_{mem}) \\
	AMAT & = 2 + 0.20 \times (15 + 0.08 \times 90)                     \\
	AMAT & = 2 + 0.20 \times (15 + 7.2)                                \\
	AMAT & = 2 + 0.20 \times 22.2                                      \\
	AMAT & = 2 + 4.44 = \mathbf{6.44 \text{ clock cycles}}
\end{align*}

\subsection{(b) AMAT with Computation Overlap and Memory Level Parallelism}
We must first calculate the effective latency for each memory hierarchy level by discounting the overlapped time and dividing by the number of simultaneous references.
\begin{align*}
	t_{L1,eff}  & = \frac{t_{L1} \times (1 - \text{Overlap}_{L1})}{\text{Refs}_{L1}} = \frac{2 \times (1 - 0.80)}{3} = \frac{0.4}{3} \approx 0.1333 \text{ cycles}  \\[10pt]
	t_{L2,eff}  & = \frac{t_{L2} \times (1 - \text{Overlap}_{L2})}{\text{Refs}_{L2}} = \frac{15 \times (1 - 0.40)}{2.5} = \frac{9.0}{2.5} = 3.6 \text{ cycles}      \\[10pt]
	t_{mem,eff} & = \frac{t_{mem} \times (1 - \text{Overlap}_{mem})}{\text{Refs}_{mem}} = \frac{90 \times (1 - 0.10)}{1.5} = \frac{81.0}{1.5} = 54.0 \text{ cycles}
\end{align*}

Now, applying these effective latencies to the AMAT formula:
\begin{align*}
	AMAT_{eff} & = t_{L1,eff} + MR_{L1} \times (t_{L2,eff} + MR_{L2} \times t_{mem,eff}) \\
	AMAT_{eff} & = 0.1333 + 0.20 \times (3.6 + 0.08 \times 54)                           \\
	AMAT_{eff} & = 0.1333 + 0.20 \times (3.6 + 4.32)                                     \\
	AMAT_{eff} & = 0.1333 + 0.20 \times 7.92                                             \\
	AMAT_{eff} & = 0.1333 + 1.584 = 1.7173 \approx \mathbf{1.72 \text{ clock cycles}}
\end{align*}

\subsection{(c) Calculate the CPI}
The perfect-cache CPI ($CPI_{perfect}$) of 0.5 is defined as having a 0-cycle access time. Therefore, we do not subtract the effective L1 hit time; the entirety of the effective AAT contributes to the final CPI.

Given that 25\% of all instructions are memory references (loads/stores):
\begin{align*}
	CPI & = CPI_{perfect} + (\text{Memory Refs per Instruction} \times AAT) \\
	CPI & = 0.5 + 0.25 \times 1.7173                                               \\
	CPI & = 0.9293 \approx \mathbf{0.93}
\end{align*}

\end{document}
